% !TeX root = proyecto.tex

%========================================================================================
% CAPÍTULO 9: CONCLUSIONES
%========================================================================================

\chapter{Conclusiones}
\label{cap:conclusiones}

El desarrollo y culminación del \textit{Sistema Inteligente para la Identificación y Seguimiento de NNA} ratifica de manera categórica e innegable el cumplimiento íntegro de los objetivos tecnológicos y sociales delineados en el protocolo de registro 2026-A135. El ciclo de vida de este proyecto atestiguó la metamorfosis de un planteamiento computacionalmente frágil, fundamentado en el análisis sintáctico lineal y el almacenamiento plano, hacia una plataforma neuro-simbólica de grado de producción, dotada de bases de datos transaccionales, infraestructura \textit{anti-bot} y visualización analítica avanzada. 

El hallazgo arquitectónico de mayor trascendencia emanado de esta investigación reside en el entendimiento empírico de las limitantes de la Inteligencia Artificial moderna frente al caos del lenguaje periodístico humano. Durante la fase experimental, quedó demostrado que depender exclusivamente del conexionismo puro de los modelos de lenguaje masivos, como los \textit{Transformers} (BETO), resulta insuficiente para un despliegue sin supervisión. Estas arquitecturas neuronales sufren de "fugas de confianza" ante narrativas semánticamente difusas, provocando alucinaciones estructurales, tales como confundir a un adolescente que comete un feminicidio con un menor de edad que ha quedado en orfandad tras el asesinato de su madre. 

La respuesta definitiva a este escollo técnico fue la conceptualización y despliegue exitoso de la Arquitectura Neuro-Simbólica. Se demostró que al injertar la estricta rigurosidad del simbolismo algorítmico tradicional (mecanismos de "Bypass" heurístico y el "Bozal de Penalización") como guardianes controladores de la inferencia neuronal, el sistema adquiere una resiliencia formidable. Esta fusión logró amalgamar la flexibilidad cognitiva de BETO con el determinismo inquebrantable de la lógica condicional, abatiendo la tasa de falsos positivos a márgenes históricamente bajos (cercanos a cero) y estableciendo un nuevo estándar metodológico para la lectura automatizada de notas fácticas de extrema sensibilidad.

Más allá de la complejidad del código y las optimizaciones matemáticas en el espacio hiperdimensional de BERTopic, el éxito supremo de este Trabajo Terminal II se materializa en su impacto social directo. La solución entregada trasciende el entorno académico para convertirse en una herramienta de alivio operativo real. Organizaciones civiles y colectivos defensores de los derechos humanos, como la fundación \textit{Futuro con Derechos}, quedan liberadas del agotador desgaste psicológico y logístico que impone la curaduría manual de la prensa roja. Al automatizar la recolección, deduplicación y clasificación de los eventos de orfandad, el sistema reduce intervenciones manuales de horas o días a escasos segundos, proveyendo bases de datos estructuradas que habilitan una acción rápida, priorizada y estratégicamente focalizada sobre las víctimas.

Finalmente, este proyecto reafirma una de las máximas más nobles de la ingeniería de sistemas: el poder computacional carece de valor si no está anclado en el profundo respeto por la dignidad humana. A pesar de procesar e ingerir volúmenes masivos de datos periodísticos a velocidades maquinales, contabilizar tragedias y clústeres geográficos, el sistema se diseñó bajo una inquebrantable doctrina de privacidad \textit{by design}. La arquitectura clasifica, cuantifica y visibiliza la urgencia de los menores, pero jamás expone sus nombres ni vulnera su derecho inalienable a la privacidad, operando en absoluta consonancia con los más altos estándares de la ética informática y la protección integral de las infancias vulneradas en México.
